\documentclass{article}
\usepackage{lmodern}
\usepackage{amsmath}
\usepackage{listings}
\usepackage{fullpage}
\usepackage{parskip}
\usepackage{xcolor}
\lstset{
	basicstyle=\ttfamily,
	columns=fullflexible,
	frame=single,
	breaklines=true,
	postbreak=\mbox{\textcolor{red}{$\hookrightarrow$}\space},
}
\begin{document}
\title{Experiment `p\_4\_3\_d\_count\_vowels' Results}
\author{\today}
\date{}
\maketitle
\textbf{Experiment outcome: }SUCCESS
\\\textbf{Bad responses: }0
\\\textbf{Responses containing}~\texttt{assume}~\textbf{: }0
\\\textbf{Responses containing}~\texttt{expect}~\textbf{: }0
\\\textbf{Resolution attempts: }1
\\\textbf{Hard fails (resolution): }0
\\\textbf{Soft fails (resolution): }0
\\\textbf{Verification attempts: }1
\section*{Problem Specification}
\textbf{Problem name: }p\_4\_3\_d\_count\_vowels
\\\textbf{Natural language statement: }null
\\\textbf{Method signature: }p\_4\_3\_d\_count\_vowels(s: string) returns (count: int)
\subsection*{Ensures}
\begin{itemize}
\item \texttt{count == count\_vowels(s) }
\end{itemize}
\subsection*{Functional Code Given}
\begin{lstlisting}
function count_vowels(s: string): int {
  if s == "" then 0
  else if s[0] in ['a','e','i','o','u','A','E','I','O','U'] then count_vowels(s[1..]) + 1
  else count_vowels(s[1..])
}
\end{lstlisting}
\clearpage
\section*{GenAI interactions}
Below you will find all interactions between the `user' (program) and the `assistant' (OpenAI).
\subsection*{Program $\rightarrow$ GenAI}
\begin{lstlisting}
You are given the following task to perform in Dafny:

Write a method to satisfy the below requirements.

The signature should be:

p_4_3_d_count_vowels(s: string) returns (count: int)

The method should respect the following contract:

ensures count == count_vowels(s) 

The contract uses the following dafny code:

function count_vowels(s: string): int {
  if s == "" then 0
  else if s[0] in ['a','e','i','o','u','A','E','I','O','U'] then count_vowels(s[1..]) + 1
  else count_vowels(s[1..])
}

These function(s) must not be used in your implementation of the method.

Produce and show only the Dafny body of this method, including the curly braces that surround it. Do not show the signature nor contract. You must not use 'assume' nor 'expect' anywhere in your code.
\end{lstlisting}
\subsection*{GenAI $\rightarrow$ Program}
\textbf{System fingerprint: }null
\\\textbf{ID: }chatcmpl-DnLsPEaCc8puF001rzt9bD03sRAaN
\\\textbf{Model: }gpt-5.5-2026-04-23
\\\textbf{Created at: }1780654537
\\\textbf{Finish reason: }stop
\\\textbf{Usage: }444 tokens in, and 1991 tokens out
\begin{lstlisting}
{
  if s == "" {
    count := 0;
  } else {
    var rest := p_4_3_d_count_vowels(s[1..]);
    if s[0] in ['a', 'e', 'i', 'o', 'u', 'A', 'E', 'I', 'O', 'U'] {
      count := rest + 1;
    } else {
      count := rest;
    }
  }
}
\end{lstlisting}
\section*{Final Program}
\begin{lstlisting}
function count_vowels(s: string): int {
  if s == "" then 0
  else if s[0] in ['a','e','i','o','u','A','E','I','O','U'] then count_vowels(s[1..]) + 1
  else count_vowels(s[1..])
}

method p_4_3_d_count_vowels(s: string) returns (count: int)
	ensures count == count_vowels(s) 
{
  if s == "" {
    count := 0;
  } else {
    var rest := p_4_3_d_count_vowels(s[1..]);
    if s[0] in ['a', 'e', 'i', 'o', 'u', 'A', 'E', 'I', 'O', 'U'] {
      count := rest + 1;
    } else {
      count := rest;
    }
  }
}
\end{lstlisting}
\section*{Total Token Usage}
\textbf{Input tokens: }444
\\\textbf{Output tokens: }1991
\\\textbf{Reasoning tokens: }1847
\\\textbf{Sum of `total tokens': }2435
\section*{Experiment Timings}
\textbf{Overall Experiment} started at 1780654536650, ended at 1780654605840, lasting 69190ms (69.19 seconds)
\\\textbf{Iteration \#1} started at 1780654536669, ended at 1780654605840, lasting 69171ms (69.17 seconds)
\end{document}
